% theme: organisms_and_environment
\MajorQuestion{生物と環境}
\SetRowSpaceQanda

\Instruction{自然界の生物どうしの関係に関する次の問いに答えなさい。}
\QQ{生物と、それを取り巻く環境を一つのまとまりとしてとらえたものを\NamedBlank{ecosystem}という。}{}
\QQ{生物どうしの食べる・食べられるという関係が、鎖のようにつながったものを\NamedBlank{food_chain}という。}{}
\QQ{実際の自然界で、複数の\RefBlank{food_chain}が網の目のように複雑につながったものを何というか。}{}
\QQ{光合成によって、無機物から有機物をつくり出す生物を\NamedBlank{producer}という。}{}
\QQ{\RefBlank{producer}がつくった有機物を、直接または間接的に取り入れる生物を何というか。}{}

\Instruction{生物の数量関係と物質の移動に関する次の問いに答えなさい。}
\QQ{食べる側ほど個体数が少なくなる数量関係を、段階ごとに積み重ねて表した形を何というか。}{}
\QQ{ある物質の生物体内での濃度が外界より高くなり、食べる・食べられる関係を通してさらに高まる現象を\NamedBlank{bioaccumulation}という。}{}
\QQ{生物の死がいや排出物などの有機物を、無機物に分解する過程に関わる生物を\NamedBlank{decomposer}という。}{}
\QQ{カビやキノコのなかまを\NamedBlank{fungi}という。}{}
\QQ{\RefBlank{fungi}の体をつくる、糸状の構造を何というか。}{}

\Instruction{土の中の微生物の特徴に関する次の問いに答えなさい。}
\QQ{\RefBlank{fungi}がつくり、なかまをふやすもとになるものを何というか。}{}
\QQ{乳酸菌や大腸菌などのなかまを\NamedBlank{bacteria}という。}{}
\QQ{一つの細胞だけで体ができている生物を何というか。}{}
\QQ{\RefBlank{bacteria}が一つの体を二つに分けてふえることを何というか。}{}
\QQ{炭素や酸素などが形を変えながら\RefBlank{ecosystem}の中をめぐることを何というか。}{}

\Instruction{自然環境の調査と人間の活動に関する次の問いに答えなさい。}
\QQ{自然環境を積極的に維持しようとすることを\NamedBlank{conservation}という。}{}
\QQ{人間の活動によって、ほかの地域から持ち込まれた生物を何というか。}{}
\QQ{ある地域にもともと定着していた生物を何というか。}{}
\QQ{地球全体の平均気温が長期的に上昇する現象を何というか。}{}
\QQ{個体数が減少し、絶滅するおそれがある生物の種を何というか。}{}
